% ============================================================
% pathA_ladder_coil_circuit.tex — coil circuit, HARDWARE ON HAND.
%
% The Path A retrofit as an electrician sees it:
%
%   L1 -> RX -> STOP -> [START || M1 aux] -> OL -> M1 -> L2
%
% RX is the VONVOFF B07CTL3TG6 receiver, programmed to MOMENTARY
% mode, held closed by the ESP32's 433 MHz heartbeat.
%
% Two things make this sheet different from ladder_coil_circuit.svg
% (the Path B sheet), and both are consequences of the actual part:
%
%  1. The receiver is a LINE-POWERED device, not a dry relay module.
%     It needs its own AC supply across the control rails, and that
%     supply must be tapped UPSTREAM of the seal-in network. Tap it
%     downstream and the receiver is dead until START is pressed,
%     its contact is open while it boots, and the seal-in can never
%     latch — the saw simply never starts.
%
%  2. Its contact therefore sits at the HEAD of the rung rather
%     than between the seal-in and the coil. Head-of-rung is the
%     placement that works whether the OUT pair turns out to be a
%     dry contact or an internally-derived switched line — see the
%     "verify the output type" legend.
%
% There is no TSTAT on this sheet. The bimetallic thermostat is
% Path B (TASK-6) and has not been purchased.
%
% Regenerate with:  make pathA_ladder_coil_circuit.svg
% ============================================================
\def\pgfsysdriver{pgfsys-dvisvgm.def}
\documentclass[border=10pt,dvisvgm]{standalone}
% ============================================================
% tsstyle.tex — shared preamble for every sheet in this directory.
%
% Inputted by each sheet after \documentclass. Sheets stay pure
% geometry; everything about how the geometry *looks* lives here,
% so a style change is a one-file edit.
% ============================================================

\usepackage[T1]{fontenc}
\usepackage[utf8]{inputenc}
\usepackage[scaled=0.92]{helvet}
\renewcommand{\familydefault}{\sfdefault}

\usepackage[american,siunitx,RPvoltages]{circuitikz}
\usetikzlibrary{calc,positioning,fit,backgrounds,arrows.meta,decorations.pathreplacing}

% ------------------------------------------------------------
% Palette. Deliberately small: ink for conductors, mute for
% annotation, and exactly two semantic colours — mains (danger)
% and SELV (safe). Colour carries meaning on these sheets; it is
% never decoration.
% ------------------------------------------------------------
\definecolor{tsink}{HTML}{16181D}    % conductors, symbol bodies
\definecolor{tsmute}{HTML}{6E7178}   % secondary annotation
\definecolor{tsrule}{HTML}{C7CAD1}   % boxes, leader lines
\definecolor{tsmains}{HTML}{B3261E}  % mains potential — danger
\definecolor{tsselv}{HTML}{15607A}   % SELV / low-voltage rail
\definecolor{tssafe}{HTML}{1E6B45}   % fail-safe / passive layer
\definecolor{tspanel}{HTML}{F5F6F7}  % block fills
\definecolor{tspanelb}{HTML}{E4E7EB} % block fills, alternate

% ------------------------------------------------------------
% Global circuitikz sizing. Bipoles are scaled down from the
% default because these sheets carry a lot of elements per unit
% area; the default American resistor is comically large next to
% an IC block.
% ------------------------------------------------------------
\ctikzset{
  bipoles/length=0.85cm,
  resistors/scale=0.85,
  capacitors/scale=0.8,
  label/align=straight,
  font=\footnotesize,
}

\tikzset{
  % Conductors -------------------------------------------------
  wire/.style        = {tsink, line width=0.9pt, line cap=round},
  buswire/.style     = {tsink, line width=1.5pt, line cap=round},
  mainswire/.style   = {tsmains, line width=1.5pt, line cap=round},
  selvwire/.style    = {tsselv, line width=1.2pt, line cap=round},
  % Functional blocks -----------------------------------------
  block/.style       = {draw=tsink, line width=1pt, fill=tspanel,
                        rounded corners=1.5pt, align=center,
                        inner sep=6pt, font=\small\bfseries},
  subblock/.style    = {draw=tsrule, line width=0.7pt, fill=white,
                        rounded corners=1.5pt, align=center,
                        inner sep=4pt, font=\scriptsize},
  % Text ------------------------------------------------------
  pinlbl/.style      = {font=\scriptsize, tsink, inner sep=2pt},
  siglbl/.style      = {font=\scriptsize\itshape, tsmute, inner sep=1.5pt},
  netlbl/.style      = {font=\scriptsize\bfseries, tsink, inner sep=2pt},
  note/.style        = {font=\scriptsize, tsmute, align=left},
  danger/.style      = {font=\scriptsize\bfseries, tsmains, align=left},
  % Isolation boundary ----------------------------------------
  isoline/.style     = {tsmains, line width=1pt, dash pattern=on 5pt off 3pt},
  % A wire that crosses another without connecting. The white
  % under-stroke reads as an over-pass, so a crossing can never be
  % mistaken for a junction (junctions are always a filled dot).
  overpass/.style    = {preaction={draw, white, line width=3.4pt}},
  % Sheet furniture -------------------------------------------
  titleblock/.style  = {draw=tsink, line width=1pt, fill=white,
                        align=left, inner sep=7pt, font=\scriptsize},
  legendbox/.style   = {draw=tsrule, line width=0.7pt, fill=white,
                        align=left, inner sep=7pt, font=\scriptsize},
}

% ------------------------------------------------------------
% \pinrow{<x>}{<y>}{<dir>}{<label>}  — a pin stub on a block
% edge. dir is -1 (left edge, stub runs left) or +1 (right edge).
% ------------------------------------------------------------
\newcommand{\pinstub}[5]{%
  \draw[wire] (#1,#2) -- ++(#3*0.55,0);
  \node[pinlbl, anchor=#4] at (#1+#3*0.10,#2+0.14) {#5};
}

% ------------------------------------------------------------
% \titleblocknode{<x>}{<y>}{<sheet>}{<subtitle>}{<source>}
% Standard title block, anchored at its south east corner.
% ------------------------------------------------------------
\newcommand{\sheettitle}[5]{%
  \node[titleblock, anchor=south east] at (#1,#2) {%
    {\normalsize\bfseries #3}\\[2pt]
    {\color{tsmute}#4}\\[4pt]
    {\color{tsmute}Table-saw thermal protection retrofit\hfill}\\
    {\color{tsmute}Generated from \texttt{#5} — do not edit the SVG}%
  };
}


% ------------------------------------------------------------
% NEMA / IEC ladder primitives, same half-gap as the Path B sheet
% so the two drawings read as one set.
% ------------------------------------------------------------
\def\cgap{0.16}
\def\cbar{0.34}

% \noc{<cx>}{<cy>}  — normally-open contact   -| |-
\newcommand{\noc}[2]{%
  \draw[wire] (#1-\cgap,#2-\cbar) -- (#1-\cgap,#2+\cbar);
  \draw[wire] (#1+\cgap,#2-\cbar) -- (#1+\cgap,#2+\cbar);
}
% \ncc{<cx>}{<cy>}  — normally-closed contact -|/|-
\newcommand{\ncc}[2]{%
  \noc{#1}{#2}%
  \draw[wire] (#1-0.34,#2-0.46) -- (#1+0.34,#2+0.46);
}
% \coilsym{<cx>}{<cy>} — output coil  -( )-
\newcommand{\coilsym}[2]{%
  \draw[wire] (#1,#2) circle (0.36);
}
% \dlabel{<cx>}{<name>}{<detail>} — name above, detail below
\newcommand{\dlabel}[3]{%
  \node[font=\small\bfseries, tsink, anchor=south] at (#1,0.86) {#2};
  \node[note, anchor=north, align=center] at (#1,-0.86) {#3};
}

\begin{document}
\begin{circuitikz}[line width=0.9pt]

% ============================================================
% Provenance banner
% ============================================================
\node[font=\scriptsize\bfseries, tssafe, anchor=south west] at (0,8.4)
  {PATH A --- DRAWN FOR THE HARDWARE ON HAND (BUILD-TONIGHT.md).
   The only added part is \texttt{B07CTL3TG6}.};

% ============================================================
% Retrofit highlight — drawn first, behind everything
% ============================================================
\begin{scope}[on background layer]
  \fill[tssafe!7, rounded corners=4pt] (1.0,-1.0) rectangle (7.4,5.6);
  \draw[tssafe!35, line width=0.8pt, rounded corners=4pt]
        (1.0,-1.0) rectangle (7.4,5.6);
\end{scope}
\node[font=\scriptsize\bfseries, tssafe, anchor=north] at (4.2,-1.15)
  {ADDED BY THE RETROFIT};

% ============================================================
% Rails
% ============================================================
\draw[buswire] (0,7.6) -- (0,-5.0);
\draw[buswire] (26,7.6) -- (26,-5.0);
\node[font=\small\bfseries, anchor=south] at (0,7.7) {L1};
\node[font=\small\bfseries, anchor=south] at (26,7.7) {L2};
\node[note, anchor=north west, align=left] at (0.15,7.62)
  {control circuit hot};
\node[note, anchor=north east, align=right] at (25.85,7.62)
  {control circuit return};
\node[note, anchor=east] at (-0.25,0) {Rung 1};

% ============================================================
% Rung — wire segments between elements
% ============================================================
\draw[wire] (0,0)     -- (2.6,0);
\draw[wire] (5.4,0)   -- (8.04,0);
\draw[wire] (8.36,0)  -- (10.4,0);
\draw[wire] (10.4,0)  -- (12.54,0);
\draw[wire] (12.86,0) -- (15.0,0);
\draw[wire] (15.0,0)  -- (20.44,0);
\draw[wire] (20.76,0) -- (23.04,0);
\draw[wire] (23.76,0) -- (26,0);

% ============================================================
% VONVOFF receiver — line-powered, four terminals
% ============================================================
\draw[block, fill=white] (1.4,2.4) rectangle (7.0,5.2);
\node[font=\small\bfseries, tsink, anchor=north] at (4.2,5.08)
  {VONVOFF RF receiver};
\node[note, anchor=north, align=center] at (4.2,4.58)
  {\texttt{B07CTL3TG6} --- 433\,MHz, AC 100--240\,V,\\
   30\,A rated on a 40\,A relay, 328\,ft.\\
   \textbf{Programmed to MOMENTARY} (1 learn press).};
\node[note, anchor=north, align=center, tssafe] at (4.2,3.28)
  {Closed only while the ESP32's carrier is arriving.};

\node[pinlbl, anchor=east] at (2.45,5.62) {AC-IN L};
\node[pinlbl, anchor=west] at (5.55,5.62) {AC-IN N};
\node[pinlbl, anchor=east] at (2.45,2.1) {OUT 1};
\node[pinlbl, anchor=west] at (5.55,2.1) {OUT 2};

% ---- its own supply, tapped ACROSS the rails, upstream of everything ----
\draw[wire] (2.6,5.2) -- (2.6,6.4);
\draw[wire] (2.6,6.4) -- (0,6.4);
\draw (0,6.4) node[circ]{};
\draw[wire] (5.4,5.2) -- (5.4,7.0);
\draw[wire] (5.4,7.0) -- (26,7.0);
\draw (26,7.0) node[circ]{};
\node[note, anchor=north west, align=left, tssafe] at (8.2,5.6)
  {\textbf{The receiver's own supply, tapped across the rails ahead of
   everything else.}\\
   It has to be: this is a line-powered unit, not a dry relay module.
   Feed it from\\
   downstream of the seal-in and it is unpowered --- and its contact
   open --- at the\\
   moment START is pressed, so the coil never latches and the saw
   never runs.};

% ---- the rung passes up through its output pair ----
\draw[wire] (2.6,0) -- (2.6,2.4);
\draw[wire] (5.4,2.4) -- (5.4,0);

% ============================================================
% Elements, left to right
% ============================================================
\ncc{8.2}{0}
\dlabel{8.2}{STOP}{NC pushbutton}

\noc{12.7}{0}
\dlabel{12.7}{START}{NO pushbutton,\\momentary}

\noc{12.7}{-3.2}
\node[font=\small\bfseries, tsink, anchor=south] at (12.7,-2.34) {M1 aux};
\node[note, anchor=north, align=center] at (12.7,-4.06)
  {NO auxiliary contact on M1.\\Carries the seal after START releases.};

\ncc{20.6}{0}
\dlabel{20.6}{OL}{NC overload aux}

\coilsym{23.4}{0}
\dlabel{23.4}{M1}{contactor coil,\\Gould/ITE A202C,\\NEMA Size 1}

% ============================================================
% Seal-in branch — M1 aux in parallel with START
% ============================================================
\draw[wire] (10.4,0) -- (10.4,-3.2);
\draw[wire] (15.0,0) -- (15.0,-3.2);
\draw[wire] (10.4,-3.2) -- (12.54,-3.2);
\draw[wire] (12.86,-3.2) -- (15.0,-3.2);
\draw (10.4,0) node[circ]{};
\draw (15.0,0) node[circ]{};

% ============================================================
% Ghost slot — where Path B's thermostat lands
% ============================================================
\draw[tsrule, line width=0.7pt, densely dashed, rounded corners=3pt]
      (16.1,1.05) rectangle (19.6,2.35);
\node[note, anchor=north, align=center, tsmute] at (17.85,2.25)
  {TSTAT lands here in Path B\\(TASK-6, not purchased).\\
   Today this span is plain wire.};
\draw[tsrule, line width=0.7pt, densely dashed, -{Latex[length=4pt]}]
      (17.85,1.05) -- (17.85,0.18);

% ============================================================
% Legends
% ============================================================
\node[legendbox, anchor=north west, text width=10.4cm] at (0,-5.8) {%
  \textbf{Before you wire it --- verify what the OUT pair actually is}\\[4pt]
  These receivers are sold two ways and the listing does not say
  which you have. With the unit \emph{unpowered}, meter continuity
  from \texttt{AC-IN L} to \texttt{OUT 1} and to \texttt{OUT 2}:\\[4pt]
  \begin{tabular}{@{}p{3.3cm}@{\hspace{5pt}}p{6.0cm}@{}}
    \textbf{Meter says} & \textbf{What you have}\\[2pt]
    \texttt{OUT} isolated from \texttt{AC-IN} on both pins
      & A genuine dry contact. It may sit anywhere in the rung,
        but leave it where this sheet draws it.\\[3pt]
    \texttt{OUT 1} bonded to \texttt{AC-IN L} (or to \texttt{N})
      & An internally-derived switched line. Head-of-rung is then
        the \emph{only} correct placement.\\
  \end{tabular}\\[5pt]
  The placement drawn here is correct in both cases, which is why
  it is drawn this way. Do not move it to sit between the seal-in
  and the coil --- that is the Path B position and it assumes a dry
  relay module the project does not own.
};

\node[legendbox, anchor=north west, text width=8.6cm] at (11.0,-5.8) {%
  \textbf{Programming MOMENTARY --- the one step that matters}\\[4pt]
  Get this wrong and the system is fail-\emph{dangerous}: a latched
  relay holds its last state when the link dies.\\[4pt]
  \begin{tabular}{@{}p{1.6cm}@{\hspace{5pt}}p{6.0cm}@{}}
    8 presses & clears every paired code --- start here\\[2pt]
    \textbf{1 press} & \textbf{point movement (momentary)} --- required\\[2pt]
    2 presses & self-locking (toggle) --- fail-dangerous\\[2pt]
    3 presses & interlock --- fail-dangerous, and the factory default\\
  \end{tabular}\\[5pt]
  Then pair a fob, \textbf{hold} its ON button, and confirm the
  contact drops within about a second of release. Time that decay
  with a stopwatch and record it --- it is the budget the heartbeat
  has to beat, and § 7 step 1 tests against it.
};

\node[legendbox, anchor=north west, text width=10.4cm] at (0,-9.9) {%
  \textbf{What opens the rung}\\[4pt]
  Every element is either \emph{NC that opens on fault} or \emph{NO
  that must be actively held closed}. No failure leaves the coil
  energised.\\[5pt]
  \begin{tabular}{@{}p{1.35cm}@{\hspace{5pt}}p{1.15cm}@{\hspace{5pt}}p{5.9cm}@{}}
    \textbf{Element} & \textbf{Type} & \textbf{Opens on}\\[2pt]
    RX     & NO & loss of the 433\,MHz carrier --- ESP32 power loss,
                  watchdog reset, NTC fault, overtemp, dead fob cell,
                  jamming, or GPIO26 floating\\[1pt]
    STOP   & NC & operator command; a broken wire also opens\\[1pt]
    START  & NO & --- (must be held until M1 aux seals in)\\[1pt]
    M1 aux & NO & loss of coil power breaks the seal\\[1pt]
    OL     & NC & sustained overcurrent\\
  \end{tabular}
};

\node[legendbox, anchor=north west, text width=8.6cm] at (11.0,-9.9) {%
  \textbf{What Path A does not have}\\[4pt]
  \textcolor{tsmains}{\textbf{No passive layer.}} SR-3 is unmet. The
  bimetallic thermostat is Path B, and until it is fitted every
  protective function on this rung depends on the ESP32 continuing
  to transmit.\\[4pt]
  \textcolor{tsmains}{\textbf{Frame, not winding.}} The DROK NTC reads
  the motor frame. It trails winding temperature, so thresholds are
  set from an observed baseline (§ 7 steps 8--9), not from the
  110\,\textdegree C figure used in ARCHITECTURE.md.\\[4pt]
  \textcolor{tssafe}{\textbf{SR-4 is met.}} Once the coil drops the
  seal-in breaks, so a human presses START. The receiver re-closing
  on its own never restarts the saw.
};

% ============================================================
% Title block + danger banner
% ============================================================
\node[titleblock, anchor=north west, text width=6.6cm] at (0,-15.4) {%
  {\normalsize\bfseries Coil circuit --- Path A (parts on hand)}\\[2pt]
  {\color{tsmute}3-wire seal-in, RF receiver at the head}\\[5pt]
  {\color{tsmute}Table-saw thermal protection retrofit}\\
  {\color{tsmute}Board side: \texttt{pathA\_supervisor.svg}}\\
  {\color{tsmute}Source: \texttt{pathA\_ladder\_coil\_circuit.tex}}%
};

\node[danger, anchor=north west, text width=11.0cm] at (11.0,-15.4) {%
  Every conductor on this sheet is at mains potential, including both
  receiver supply taps. Lock out the disconnect before touching the
  starter enclosure, and verify 0\,V with a meter you have just proved
  on a known-live circuit.};

\end{circuitikz}
\end{document}
